% SGA 3，第 VII_A 讲， 英文重建稿——第 1 轮。
% 范围：权威底本局部页 52--55, 节 8 开头至 8.3.1,
% 止于其前： 8.4 局部页 56.
% 合并阅读本页码： 505--508. 再版印刷页码： 494--497.
% 本范围内原版页边定位： 470--472.
% 权威底本： Polo--Gille Exp7A-13oct24.pdf,
% SHA-256 D46DB8C79D0385643B6EED4B96900613248375F9B73B4202478A4F57FD3E8522.
% 曾参考 Jacob Reinhold 的 CC BY 4.0 Markdown 作为比对谱系。

\section*{8. 基底为域的情形}
\phantomsection
\label{sga3:VIIA:8}
\addcontentsline{toc}{section}{8. 基底为域的情形}

\medskip
\noindent\textbf{8.1.}
\phantomsection
\label{sga3:VIIA:8.1}

% 权威底本： Exp7A 局部页 52; 合并本页 505; 印刷页 494.
现在汇集当 \(S\) 是特征为 \(p>0\) 的域 \(k\) 的谱时所得的结果。若一个 \(k\)-群概形的底层概形在 \(k\) 上为有限型，则称它是\emph{代数的}。\SGARefLink{sga3:VIIA:7.2}{定理 7.2} 给出下述结果。\footnote{编者添加了从 \SGARefLink{sga3:VIIA:8.1.1}{8.1.1} 到 \SGARefLink{sga3:VIIA:8.1.3}{8.1.3} 的编号，以突出这里陈述的结果。}

\medskip
\noindent\textbf{定理 8.1.1。}
\phantomsection
\label{sga3:VIIA:8.1.1}
\textit{函子 \(\mathfrak G_p\) 将每个 \(k\) 上的有限维受限 \(p\)-Lie 代数 \(\mathcal L\) 对应到 \(k\)-群 \(\mathfrak G_p(\mathcal L)\)，它是将每个代数 \(k\)-群 \(G\) 对应到 \(\operatorname{Lie}(G)\) 的函子的左伴随。}

将此结果与 \SGARefLink{sga3:VIIA:7.4}{定理 7.4} 结合，得到：

\medskip
\noindent\textbf{定理 8.1.2。}
\phantomsection
\label{sga3:VIIA:8.1.2}
\textit{函子
\[
  \mathfrak G_p:
  \mathcal L\longmapsto\mathfrak G_p(\mathcal L),
  \qquad
  G\longmapsto\operatorname{Lie}(G),
\]
在 \(k\) 上有限维受限 \(p\)-Lie 代数的范畴与高度至多为 \(1\) 的代数 \(k\)-群的范畴之间诱导出互为拟逆的范畴等价。}

% 权威底本： Exp7A 局部页 53; 合并本页 506; 印刷页 495.
由于 \(\mathfrak G_p\) 是左伴随，它与归纳极限交换，\footnote{编者扩充了下面的论证以及 \SGARefLink{sga3:VIIA:8.1.3}{推论 8.1.3} 的证明。}特别地，它与余核交换。设
\[
  \varphi:G\longrightarrow H,
  \qquad
  \varphi':G'\longrightarrow H
\]
是高度至多为 \(1\) 的代数 \(k\)-群之间的两个态射。纤维积 \(G\times_HG'\) 仍是高度至多为 \(1\) 的代数 \(k\)-群，因为 Frobenius 态射
\[
  \operatorname{Fr}:G\longrightarrow G^{(p)}
\]
与纤维积交换。因此，将高度至多为 \(1\) 的代数 \(k\)-群的范畴包含到所有代数 \(k\)-群的范畴中的函子与纤维积交换，因而特别与核交换。得到：

\medskip
\noindent\textbf{推论 8.1.3。}
\phantomsection
\label{sga3:VIIA:8.1.3}
\textit{函子 \(\mathfrak G_p\) 在下述意义下是正合的。设
\[
  \pi:\mathcal L_1\longrightarrow\mathcal L_2
\]
是 \(k\) 上有限维受限 \(p\)-Lie 代数之间的满同态，并令 \(i\) 为 \(\mathcal L_0=\ker\pi\) 到 \(\mathcal L_1\) 的包含映射。则下列代数 \(k\)-群的列正合：}
\begin{center}
  \phantomsection
  \label{sga3:VIIA:8.1.3-Gp-exact-sequence}
  % SGA3-VII-DIAG-071
  \(\displaystyle
1 \longrightarrow \mathfrak G_p(\mathcal L_0)
  \xrightarrow{\mathfrak G_p(i)}
\mathfrak G_p(\mathcal L_1)
  \xrightarrow{\mathfrak G_p(\pi)}
\mathfrak G_p(\mathcal L_2)
  \longrightarrow 1 .
\)

\end{center}
\footnote{此外，由 VIA，\SGARefLink{thm:sga3-VIA-3.2}{3.2}，\(\mathfrak G_p(\mathcal L_2)\) 表示 \(\mathfrak G_p(\mathcal L_1)\) 对 \(\mathfrak G_p(\mathcal L_0)\) 的 (fppf) 商层。}

事实上，由上述讨论，\(\mathfrak G_p(i)\) 诱导出 \(\mathfrak G_p(\mathcal L_0)\) 到 \(\ker(\mathfrak G_p(\pi))\) 的同构。不论在所有代数 \(k\)-群的范畴中取核，还是在高度至多为 \(1\) 的代数 \(k\)-群的子范畴中取核，所得的核都相同。此外，
\[
  \mathfrak G_p(\pi):
  \mathfrak G_p(\mathcal L_1)
  \longrightarrow
  \mathfrak G_p(\mathcal L_2)
\]
将 \(\mathfrak G_p(\mathcal L_2)\) 等同于代数 \(k\)-群范畴中的 \(\mathfrak G_p(\mathcal L_1)\) 对 \(\mathfrak G_p(\mathcal L_0)\) 的商。

\medskip
\noindent\textbf{注 8.1.4。}\footnote{编者添加了由 O.~Gabber 告知的此注，以便在 \SGARefLink{sga3:VIIA:8.3.1}{8.3.1} 中使用。}
\phantomsection
\label{sga3:VIIA:8.1.4}

设 \(\varphi:G\to H\) 是 \(k\)-群态射，并令 \(K=\ker(\varphi)\)。假设 \(\varphi\) 对 (fpqc) 拓扑是覆盖的。特别地，若它对更粗的拓扑（例如 (fppf) 拓扑）是覆盖的，就属于这种情形。此时 \(\varphi\) 是 \(H\) 上的 \(K\)-挠子（参见 IV，\SGARefLink{sga3:iv:corollary:5:1:7:1}{5.1.7.1}）。此外，由 IV，\SGARefLink{sga3:iv:proposition:6:3:1}{6.3.1}，存在由仿射开子集 \(S_i\) 构成的 \(H\) 的一个覆盖，并且对于每个 \(i\)，存在一个通过 \(\varphi\) 分解的仿射忠实平坦态射 \(T_i\to S_i\)。而 \(G\times_HT_i\) 与 \(K\times T_i\) 在 \(T_i\) 上同构，因而在 \(T_i\) 上忠实平坦。由 (fpqc) 下降，\(G\times_HS_i\to S_i\) 忠实平坦，故 \(\varphi\) 忠实平坦。

反之，若 \(\varphi\) 忠实平坦且拟紧，分别地，忠实平坦且局部有限呈示，则它对 (fpqc) 拓扑，分别地，对 (fppf) 拓扑是覆盖的（参见 IV，\SGARefLink{sga3:iv:proposition:6:3:1}{6.3.1}）。最后，层态射是覆盖的当且仅当它是满态射（参见 IV，\SGARefLink{sga3:iv:proposition:4:4:3}{4.4.3}）。特别地，一个拟紧的 \(k\)-群态射忠实平坦，当且仅当它是 (fpqc) 层的满态射。

\medskip
\noindent\textbf{8.2. 命题。}
\phantomsection
\label{sga3:VIIA:8.2}
\textit{考虑特征为 \(p>0\) 的域 \(k\) 上的代数群正合列\footnote{因此 \(\pi\) 忠实平坦，且 \(\tau\) 将 \(G'\) 等同于 \(\ker\pi\)，从而 \(G''\) 表示 \(G\) 对 \(G'\) 的 (fppf) 商层；参见 VIA，\SGARefLink{thm:sga3-VIA-3.2}{3.2} 和 \SGARefLink{par:sga3-VIA-5.2}{5.2}。}：}
\begin{center}
  \phantomsection
  \label{sga3:VIIA:8.2-algebraic-groups-exact-sequence}
  % SGA3-VII-DIAG-072
  \(\displaystyle
1 \longrightarrow G'
  \xrightarrow{\tau} G
  \xrightarrow{\pi} G''
  \longrightarrow 1 .
\)

\end{center}
\textit{考虑下列断言：
\begin{enumerate}
  \item[\textup{\SGAInlineRefTarget{sga3:VIIA:8.2:item:i}(i)}] \(\pi\) 是平滑的；
  \item[\textup{\SGAInlineRefTarget{sga3:VIIA:8.2:item:ii}(ii)}] \(G'\) 是平滑的；
  \item[\textup{\SGAInlineRefTarget{sga3:VIIA:8.2:item:iii}(iii)}] 对于每个整数 \(n>0\)，由 \(\tau\) 和 \(\pi\) 诱导的列
\end{enumerate}}
\begin{center}
  \phantomsection
  \label{sga3:VIIA:8.2-frobenius-kernel-exact-sequence}
  % SGA3-VII-DIAG-073
  \(\displaystyle
1 \longrightarrow \operatorname{Fr}^nG'
  \longrightarrow \operatorname{Fr}^nG
  \longrightarrow \operatorname{Fr}^nG''
  \longrightarrow 1 .
\)

\end{center}
\textit{\noindent 正合；
\begin{enumerate}
  \item[\textup{\SGAInlineRefTarget{sga3:VIIA:8.2:item:iv}(iv)}] 态射
  \[
    \operatorname{Fr}\pi:
    \operatorname{Fr}G\longrightarrow\operatorname{Fr}G''
  \]
  是 (fppf) 层的满态射；
  \item[\textup{\SGAInlineRefTarget{sga3:VIIA:8.2:item:v}(v)}] 态射
  \[
    \operatorname{Lie}(\pi):
    \operatorname{Lie}(G)\longrightarrow\operatorname{Lie}(G'')
  \]
  是满射。
\end{enumerate}
则
\[
  \textup{(i)}\Longleftrightarrow\textup{(ii)}
  \Longrightarrow\textup{(iii)}
  \Longrightarrow\textup{(iv)}
  \Longleftrightarrow\textup{(v)},
\]
且当 \(G\) 在 \(k\) 上平滑时，这五个断言全部等价。}

% 权威底本： Exp7A 局部页 54; 合并本页 507; 印刷页 496.
事实上，由 VIB，\SGARefLink{prop:sga3-VIB-9.2}{9.2}~\SGARefLink{concl:sga3-VIB-9.2-vii}{(vii)}，\SGARefLink{sga3:VIIA:8.2:item:i}{(i)} 等价于 \SGARefLink{sga3:VIIA:8.2:item:ii}{(ii)}，而 \SGARefLink{sga3:VIIA:8.2:item:iii}{(iii)} 显然蕴含 \SGARefLink{sga3:VIIA:8.2:item:iv}{(iv)}。\SGARefLink{sga3:VIIA:8.2:item:iv}{(iv)} 与 \SGARefLink{sga3:VIIA:8.2:item:v}{(v)} 的等价性由 \SGARefLink{sga3:VIIA:8.1.3}{8.1.3} 推出。

蕴含关系 \SGARefLink{sga3:VIIA:8.2:item:ii}{(ii)} \(\Rightarrow\) \SGARefLink{sga3:VIIA:8.2:item:iii}{(iii)} 由下图推出
\begin{center}
  \phantomsection
  \label{sga3:VIIA:8.2-frobenius-exact-rows-diagram}
  % SGA3-VII-DIAG-074
  \resizebox{0.98\linewidth}{!}{%
\begin{tikzcd}[ampersand replacement=\&,row sep=3.6em,column sep=3.5em]
1 \arrow[r]
  \& G' \arrow[r,"\tau"] \arrow[d,swap,"\operatorname{Fr}^n(G'/k)"]
  \& G \arrow[r,"\pi"] \arrow[d,swap,"\operatorname{Fr}^n(G/k)"]
  \& G'' \arrow[r] \arrow[d,swap,"\operatorname{Fr}^n(G''/k)"]
  \& 1 \\
1 \arrow[r]
  \& G'^{(p^n)} \arrow[r,"\tau^{(p^n)}"']
  \& G^{(p^n)} \arrow[r,"\pi^{(p^n)}"']
  \& G''^{(p^n)} \arrow[r]
  \& 1
\end{tikzcd}%
}

\end{center}
，其中两行均正合。由下文的 \SGARefLink{sga3:VIIA:8.3.1}{推论 8.3.1}，\(\operatorname{Fr}^n(G'/k)\) 是 (fppf) 层的满态射，因此 \(\pi\) 诱导出从 \(\operatorname{Fr}^nG\) 到 \(\operatorname{Fr}^nG''\) 的满态射；这里使用未必交换的群层的蛇引理。

同样，假设 \(G\) 在 \(k\) 上平滑。则 \(\operatorname{Fr}(G/k)\) 是满态射。若 \(\operatorname{Fr}\pi\) 也是满态射，则对上图取 \(n=1\) 并应用同一个蛇引理，可知 \(\operatorname{Fr}(G'/k)\) 是满态射。因此由下文 \SGARefLink{sga3:VIIA:8.3.1}{8.3.1}，\(G'\) 在 \(k\) 上平滑。

\medskip
\noindent\textbf{8.3. 命题。}\footnote{编者将“代数的”改为“局部有限型”。}
\phantomsection
\label{sga3:VIIA:8.3}
\textit{设 \(G\) 是特征为 \(p>0\) 的域 \(k\) 上的局部有限型群。存在整数 \(n_0\)，使得
\[
  G/(\operatorname{Fr}^nG)
\]
对于每个 \(n\ge n_0\) 都在 \(k\) 上平滑。}

\(G/(\operatorname{Fr}^nG)\) 的构造与基域扩张交换（\SGARefLink{sga3:VIIA:4.1.1}{4.1.1} 和 VIA，\SGARefLink{thm:sga3-VIA-3.3.2}{3.3.2}），因此可以假设 \(k\) 完美。于是 \(G_{\mathrm{red}}\) 是局部有限型 \(k\)-群（参见 VIA，\SGARefLink{sga3:VIA:0.2}{0.2}），并且令
\[
  H=G_{\mathrm{red}}\backslash G,
\]
，就有下述交换正合图：
\begin{center}
  \phantomsection
  \label{sga3:VIIA:8.3-reduced-quotient-frobenius-diagram}
  % SGA3-VII-DIAG-075
  \resizebox{0.94\linewidth}{!}{%
\begin{tikzcd}[ampersand replacement=\&,row sep=3.8em,column sep=3.7em]
1 \arrow[r]
  \& G_{\mathrm{red}} \arrow[r] \arrow[d,"\operatorname{Fr}^n(G_{\mathrm{red}}/k)"]
  \& G \arrow[r] \arrow[d,"\operatorname{Fr}^n(G/k)"]
  \& H \arrow[d,"\operatorname{Fr}^n(H/k)"] \\
1 \arrow[r]
  \& G_{\mathrm{red}}^{(p^n)} \arrow[r]
  \& G^{(p^n)} \arrow[r]
  \& H^{(p^n)}\, .
\end{tikzcd}%
}

\end{center}

% 权威底本： Exp7A 局部页 55; 合并本页 508; 印刷页 497.
概形 \(H\) 是一个剩余域为 \(k\) 的有限局部 \(k\)-代数的谱（参见 VIA，\SGARefLink{prop:sga3-VIA-5.6.1}{5.6.1}）。因此存在整数 \(n_0\)，使得对于每个 \(n\ge n_0\)，\(\operatorname{Fr}^n(H/k)\) 都通过 \(H^{(p^n)}\) 的单位截面分解。由此可知，对于 \(n\ge n_0\)，\(\operatorname{Fr}^n(G/k)\) 通过 \(G_{\mathrm{red}}^{(p^n)}\) 分解，而 VIA，\SGARefLink{prop:sga3-VIA-5.4.1}{5.4.1} 给出交换图
\begin{center}
  \phantomsection
  \label{sga3:VIIA:8.3-frobenius-factorization-triangle}
  % SGA3-VII-DIAG-076
  \begin{tikzcd}[row sep=3.6em,column sep=4.6em]
G \arrow[r,"\operatorname{Fr}^n(G/k)"] \arrow[d,swap,"\pi"]
  & G_{\mathrm{red}}^{(p^n)} \\
G/(\operatorname{Fr}^nG) \arrow[ur,swap,"i"]
\end{tikzcd}

\end{center}
其中 \(i\) 是闭浸入，\(\pi\) 是典范投影。此外，\(i\) 在底层拓扑空间上诱导同胚，故为同构。由于 \(k\) 完美，\(G_{\mathrm{red}}^{(p^n)}\) 在 \(k\) 上平滑（VIA，\SGARefLink{prop:sga3-VIA-1.3.1}{1.3.1}）。因此，对于每个 \(n\ge n_0\)，\(G/(\operatorname{Fr}^nG)\) 都在 \(k\) 上平滑。

\medskip
\noindent\textbf{8.3.1. 推论。}
\phantomsection
\label{sga3:VIIA:8.3.1}
\textit{设 \(G\) 是特征为 \(p>0\) 的域 \(k\) 上的局部有限型群，并设 \(n\ge1\)。\footnote{编者明确写出了 \SGARefLink{sga3:VIIA:8.3.1:item:ii}{(ii)} 与 \SGARefLink{sga3:VIIA:8.3.1:item:iii}{(iii)} 的等价性，并扩充了证明。}下列条件等价：
\begin{enumerate}
  \item[\textup{(i)}] \(G\) 在 \(k\) 上平滑；
  \item[\textup{\SGAInlineRefTarget{sga3:VIIA:8.3.1:item:ii}(ii)}]
  \(\operatorname{Fr}^n(G/k):G\to G^{(p^n)}\) 是 (fppf) 层的满态射；
  \item[\textup{\SGAInlineRefTarget{sga3:VIIA:8.3.1:item:iii}(iii)}]
  \(\operatorname{Fr}^n(G/k):G\to G^{(p^n)}\) 忠实平坦。
\end{enumerate}}

由于 \(G\) 在 \(k\) 上局部有限型，\(\operatorname{Fr}^n(G/k)\) 有限呈示，因此 \SGARefLink{sga3:VIIA:8.3.1:item:ii}{(ii)} 与 \SGARefLink{sga3:VIIA:8.3.1:item:iii}{(iii)} 的等价性由 \SGARefLink{sga3:VIIA:8.1.4}{8.1.4} 推出。假设 \(G\) 在 \(k\) 上平滑，因而既约。由于 \(\operatorname{Fr}^n(G/k)\) 是满射，它忠实平坦（参见 VIA，\SGARefLink{cor:sga3-VIA-6.2}{6.2} 或 VIB，\SGARefLink{prop:sga3-VIB-1.3}{1.3}）。

反之，假设 \(\operatorname{Fr}^n(G/k)\) 忠实平坦。态射
\[
  \operatorname{Fr}^n(G^{(p^n)}/k)
\]
由 \(\operatorname{Fr}^n(G/k)\) 通过基变换得到（参见 \SGARefLink{sga3:VIIA:4.1.3}{4.1.3}），所以它也忠实平坦，复合
\[
  \operatorname{Fr}^{2n}(G/k):
  G\longrightarrow G^{(p^n)}
  \longrightarrow G^{(p^{2n})}.
\]
也是如此。由此可知，对于每个 \(m\in\mathbb N\)，
\[
  \operatorname{Fr}^{mn}(G/k):
  G\longrightarrow G^{(p^{mn})}
\]
忠实平坦。因此它诱导出同构
\[
  G/(\operatorname{Fr}^{mn}G)
  \simeq
  G^{(p^{mn})}
\]
（参见 VIA，\SGARefLink{prop:sga3-VIA-5.4.1}{5.4.1}）。由 \SGARefLink{sga3:VIIA:8.3}{命题 8.3}，当 \(m\) 充分大时，\(G^{(p^{mn})}\) 在 \(k\) 上平滑，因此由 (fpqc) 下降，\(G\) 也平滑（参见 EGA IV\(_4\)，17.7.1）。

% 下一游标： Exp7A 局部页 56, 合并阅读本页 509,
% 再版印刷页 498, 位于 8.4.
